\begin{frame}{Proposed Method}
  \begin{columns}[T]
    \column{0.34\textwidth}
    \begin{itemize}
      \item Primary method: \textbf{MV-DUSt3R+}
      \item Single-stage feed-forward multi-view reconstruction
      \item Final setting uses sparse posed RGB-D views
      \item Run 30 adds source-depth correction after MV-DUSt3R+
      \item Our study focuses on robustness, hard cases, and runtime
    \end{itemize}

    \vspace{0.15cm}
    \figsource{Tang et al., MV-DUSt3R+, project page and paper}

    \column{0.66\textwidth}
    \centering
    \includegraphics[width=\linewidth]{mv_overview_slide.png}
  \end{columns}
\end{frame}

\begin{frame}{Architecture Comparison}
  \begin{columns}[T]
    \column{0.49\textwidth}
    \centering
    \textbf{DUSt3R baseline}\par
    \vspace{0.1cm}
    \includegraphics[width=\linewidth]{dust3r_arch_crop.png}

    \vspace{0.08cm}
    {\scriptsize Pairwise reasoning with cross-attention between two views}

    \column{0.49\textwidth}
    \centering
    \textbf{MV-DUSt3R+ main method}\par
    \vspace{0.1cm}
    \includegraphics[width=\linewidth]{mv_arch_slide.png}

    \vspace{0.08cm}
    {\scriptsize Multi-path design with cross-reference-view information flow}
  \end{columns}

  \vspace{0.12cm}
  \figsource{Left: Wang et al., DUSt3R, CVPR 2024. Right: Tang et al., MV-DUSt3R+, project page/paper}
\end{frame}

\begin{frame}{Why MV-DUSt3R+ vs DUSt3R?}
  \small
  \begin{tabular}{p{0.22\textwidth} p{0.33\textwidth} p{0.33\textwidth}}
    \toprule
    \textbf{Aspect} & \textbf{DUSt3R} & \textbf{MV-DUSt3R+} \\
    \midrule
    Processing style & Pairwise reconstruction & Joint multi-view reconstruction \\
    Global consistency & Requires later alignment & Built into one forward pass \\
    Sparse-view behavior & Can accumulate pairwise errors & More robust cue fusion \\
    Runtime trend & Slower as views increase & Better scaling with more views \\
    Role in project & Reference baseline & Main candidate method \\
    \bottomrule
  \end{tabular}

  \vspace{0.22cm}
  \begin{itemize}
    \item DUSt3R is still a strong baseline, so the comparison is direct and meaningful.
    \item MV-DUSt3R+ is the candidate generator for both the Run 30 RGB-D
    reference and the target estimated-depth correction step.
  \end{itemize}
\end{frame}

\begin{frame}{Quantitative Evidence from the Paper}
  \centering
  \includegraphics[height=0.62\textheight,keepaspectratio]{mvdust3r_results_table.png}

  \vspace{0.04cm}
  {\scriptsize Across multiple datasets and view counts, the highlighted rows show that MV-DUSt3R+ improves reconstruction quality while keeping runtime practical.}

  \figsource{Tang et al., MV-DUSt3R+, Table 2}
\end{frame}

\begin{frame}{Qualitative Motivation}
  \centering
  \includegraphics[height=0.65\textheight,keepaspectratio]{mv_qualitative_slide.png}

  \vspace{0.04cm}
  {\scriptsize DUSt3R often fails on repeated structures, while MV-DUSt3R+ remains more stable in sparse-view scenes.}

  \figsource{Tang et al., MV-DUSt3R+, qualitative results}
\end{frame}
